\section{Introduction}
\label{sec:intro}
The rapid development of information technology has resulted in an enormous amount of data, posing significant challenges to data management systems \cite{bigdata16}. Indexes are vital components that enhance query performance by avoiding exhaustive searches. Traditional indexes \cite{btree06, btree11, bwtree13, dynamichash98, hash16, lsmtree96, lsmtree17, remix21, adaptive18, adaptive22}, such as B+Tree, do not take full advantage of data distributions, leading to inefficiencies caused by pointer chasing and cache misses \cite{learned23}. Learned indexes \cite{case18}, proposed in recent years, regard index as a model of data distribution, and makes proper use of this model to improve query efficiency and reduce index size.

In recent years, learned indexes have emerged as a compelling alternative to traditional structures by treating the task of locating a record as a prediction problem. However, accurate position prediction is extremely difficult and expensive, instead, most learned indexes predict within an error threshold, then perform a bounded search around predicted position to locate the true position. To keep both training and prediction costs low, they employ simple models, such as linear model or a single-layer neural network. Moreover, to improve overall accuracy and reduce correction overhead, the data is partitioned into multiple segments, each trained with its own model so as to capture local distributions with high precision. This segmented modeling both raises prediction accuracy and enhances scalability, enabling learned indexes to handle large, complex datasets. Empirical studies \cite{sosd20} have shown that even a simple piecewise linear model works effectively on the majority of real-world datasets. Accordingly, most existing learned indexes \cite{pgm20, radixspline20, fiting-tree19, alex20, madex20, coax23, swix24, disklearned24} adopt piecewise linear regression as their prediction backbone.

Despite these gains, some evaluations \cite{sosd20, micro22, gre22, learned23} have shown that existing learned indexes still have a series of flaws. For example, \cite{case18, plex21, towards21} is only suitable for read-only workloads; \cite{xindex20, lipp21, finedex21, chameleon24} deliver excellent point queries throughput but suffer on range queries; \cite{pgm20} performs well after bulk loading but lack stability after updates; \cite{lipp21} achieves efficient queries and updates, but incur high space usage; \cite{alex20, swix24} impose slow bulk‐load times. Critically, most proposals focus on optimizing only one or two metrics at a time, making it difficult to achieve a multidimensional trade‑off that outperforms the classic B+Tree. This gap leaves learned indexes largely experimental: selecting the appropriate index for an unknown or evolving workload is itself a challenge. As a result, learned indexes struggle to gain traction in practical, real‑world settings.

To address these limitations, we propose LMIndex, a foundational structure optimized for query and space efficiency, and its variant LMG, which employs localized gap allocation to absorb dynamic updates and ultimately achieve multi-dimensional performance balance. Their design stems from the optimization of two fundamental indexing phases: (1) the segment index step, which locates the appropriate data segment. To guarantee structural robustness under dynamic workloads, this top-layer structure needs to maintain a small memory footprint, execute theoretical $O(1)$ routing, and eliminate cascading maintenance overheads. (2) the localized correction step, which searches within the target segment to find the exact record position. To sustain high-speed retrievals without sacrificing overall system balance, this process demands lower prediction errors and optimized local search algorithms. Concretely, LMIndex and LMG meticulously co-design the routing logic and local correction mechanisms to navigate inherent structural trade-offs, leading to the following contributions:

% \textbf{(1) A robust, spatio-temporal efficient segment index structure.} 
% We introduce the Linear Map (LM), an internal structure designed for indexing segments that maintains robust efficiency across dynamic workloads and diverse datasets. In practice, given a fixed key type, the LM achieves $O(1)$ complexity for both queries and updates. Empirically, it incurs an average space overhead of only 4\% in our experiments.

\textbf{(1) Robust and Efficient Segment Index Structure.} To optimize the segment index step, we introduce Linear Map (LM). Compared to other segment index structures, LM eliminates the need for complex model training. Instead, it utilizes a deterministic mapping approach that flexibly and efficiently adapts to changes in the underlying segment distribution. Consequently, given a fixed key type, LM achieves $O(1)$ time complexity for both queries and dynamic updates. Empirically, it maintains robust performance across diverse datasets while incurring an average space overhead of only 4\%.

\textbf{(2) Optimal Error Threshold Compression Algorithm for Local Correction.} To optimize the local correction step, we introduce the Optimal Error Threshold Algorithm (OETA). Compared to the least squares method, OETA minimizes the maximum prediction error within each segment while achieving performance that is independent of distribution and nearly constant in time. Furthermore, we demonstrate that binary search (BS) outperforms exponential search (ES) on complex distributions. OETA accelerates BS, enabling it to match the performance of exponential search even on highly linear distributions.

% \textbf{(3) Robust and Efficient Framework for Multi-Dimensional Balance.} By integrating the LM structure with the OETA algorithm, we first construct LMIndex, a foundational structure strictly optimized for read efficiency and compact space utilization. Building upon this baseline, we develop LMG, an write-optimized variant employing a gapped array coupled with a novel gap allocation strategy. Consequently, LMG mitigates the severe performance trade-offs inherent in state-of-the-art baselines, delivering a highly robust and efficient framework that maintains performance stability and competitive read/write throughput across all evaluation dimensions.

\textbf{(3) Robust and Efficient Framework for Multi-Dimensional Balance.} By integrating the LM structure with the OETA algorithm, we first construct LMIndex, a foundational learned index framework employing a dense data layout. To unleash its potential in dynamic environments, we further develop LMG, a variant employing a gapped data layout and a novel gap allocation strategy. LMG mitigates the performance trade-offs inherent in state-of-the-art baselines, delivering a highly robust and efficient framework that maintains performance stability and competitive read/write throughput across all evaluation dimensions.

The remainder of this paper is organized as follows. Section \ref{sec:back&moti} describes the background and motivation. Section \ref{sec:overview} provides an overview of the overall design of LMIndex and LMG, including basic operations. Section \ref{sec:lm} introduces the Linear Map in detail. Section \ref{sec:single-model} introduces Optimal Error Threshold Algorithm. Section \ref{sec:evaluation} presents experimental analysis. Section \ref{sec:extend} explores the extensibility of our framework to broader database operations. Section \ref{sec:rel-work} reviews related work in the field. Finally, Section \ref{sec:conclusions} concludes the paper.